\documentclass[11pt]{letter}

\usepackage{geometry}
\geometry{margin=1in}
\usepackage{setspace}
\usepackage{hyperref}

\signature{Decory Edwards}
\address{Decory Edwards \\ Johns Hopkins University \\ Department of Economics}

\begin{document}

\begin{letter}{Hiring Committee \\ Applied Academics \\ New York, NY}

\opening{Dear Hiring Committee,}

I am writing to apply for the Client Strategy \& Solutions Manager position at Applied Academics. I am currently a PhD candidate in Economics at Johns Hopkins University. My training combines rigorous quantitative modeling, applied data analysis, and careful economic reasoning. I am motivated by work that sits at the intersection of research, investment decision making, and real-world implementation. Applied Academics’ model of pairing first-principles quantitative research with deep client engagement strongly aligns with how I think about impactful economic work.

My research agenda is at the intersection of wealth inequality and macroeconomics. In my job market paper, I build and estimate a structural heterogeneous agent model with returns that differ across households. The core contribution is to show how heterogeneity in returns and income risk jointly shape the wealth distribution and consumption behavior. Relative to closely related work, my framework performs better at matching the lower and middle portions of the wealth distribution. This difference matters quantitatively. Models that focus primarily on returns heterogeneity can fit the top of the wealth distribution closely but tend to generate too much wealth among the bottom half of households. In my framework, income uncertainty generates a precautionary saving motive that produces genuinely low wealth households. This leads to a realistic distribution of marginal propensities to consume and aggregate consumption responses that align with empirical estimates.

A central feature of my work is translating complex quantitative frameworks into insights that are decision relevant. I routinely synthesize modeling results, data, and economic intuition into clear written explanations and presentations for non-technical audiences. I am comfortable working through ambiguous problems, clarifying objectives, and adapting analytical outputs to the specific constraints and priorities of stakeholders. This mirrors the challenge of helping institutional clients interpret and apply quantitative research to their portfolios.

I am particularly interested in Applied Academics because of its role as a quantitative manager of managers and its focus on institutional investors with long horizons. The opportunity to work closely with the research team while serving as a trusted partner to CIOs and investment teams is especially appealing. I value environments where intellectual rigor, clarity of communication, and long-term relationships are central to success.

My economics training has given me a strong foundation in statistical reasoning, modeling, and disciplined analysis, as well as the ability to learn complex quantitative systems quickly. I bring a strong ownership mindset, attention to detail, and a collaborative approach to problem solving. I enjoy taking responsibility for projects end to end, coordinating across teams, and ensuring high quality outcomes.

I am enthusiastic about the opportunity to contribute to Applied Academics’ growth and to help scale its impact through thoughtful client engagement and strategic problem solving. I would welcome the opportunity to bring my analytical background and communication skills to this role.

Thank you for your time and consideration.

\closing{Sincerely,}

\end{letter}
\end{document}
